\subsection{Implementação da infraestrutura}
\label{subsec:implementacao_infra}

A infraestrutura foi implementada em três componentes: o firmware do nó sensor, o firmware do \textit{gateway} e o \textit{backend} de persistência, além de ferramentas de análise. Os dois firmwares rodam em placas Heltec WiFi LoRa 32 V3 \cite{heltec_wifi_lora_32_v3}, que integram o microcontrolador ESP32-S3 e o transceptor LoRa SX1262 \cite{semtech_sx1262}. O código está em repositório público, com os registros de decisão que documentam cada escolha e as alternativas descartadas.

Os parâmetros de rádio e o formato do quadro ficam em um único cabeçalho compartilhado pelos dois firmwares. A escolha é deliberada: uma divergência de frequência, fator de espalhamento ou disposição de bytes entre nó e \textit{gateway} não produz erro de compilação nem mensagem em tempo de execução, apenas um receptor que não recebe nada. Verificações em tempo de compilação garantem que o quadro tem exatamente 26 bytes e não excede o limite de 51 bytes do plano regional.

Dois caminhos de segurança coexistem no firmware, selecionados por ambiente de compilação, com o restante do código em comum. O primeiro, de linha de base, usa AES-128-GCM com a chave de rede no cabeçalho compartilhado: o quadro leva em claro um cabeçalho de 6 bytes (versão, identificador do dispositivo e contador), autenticado como dado associado, seguido de 4 bytes cifrados e da tag de 16 bytes. O segundo, padrão desde 8 de outubro de 2026, usa o Ascon-128a pela biblioteca \textit{AsconSecure} da frente de criptografia (Subseção~\ref{subsec:asconsecure}): o cabeçalho em claro e autenticado reduz-se a 2 bytes (versão e identificador), e o pacote da biblioteca ocupa os 24 restantes, com o contador de 4 bytes, os 4 bytes cifrados e a tag de 16. Os dois quadros têm 26 bytes e o mesmo tempo no ar. No caminho Ascon a chave é de cada nó: gerada na própria placa durante o provisionamento, a partir do gerador de hardware do microcontrolador, e guardada na sua memória não volátil, sem nunca constar do código; o \textit{gateway} guarda a chave de cada nó em um arquivo fora do repositório e a escolhe pelo identificador do cabeçalho. O contador da biblioteca é o contador de quadros do registro.

O nó transmite um quadro a cada cinco segundos, com agendamento em tempo absoluto: o instante da próxima transmissão é calculado a partir do anterior, e não a partir do fim da transmissão corrente, de modo que o tempo de processamento não se acumula como atraso. O contador de quadros é persistido na memória não volátil em blocos de 32: o nó grava o teto do bloco antes de usá-lo e, ao reiniciar, retoma do próximo bloco. Alguns valores são pulados a cada reinício, mas nenhum se repete, o que preserva a unicidade do \textit{nonce} derivado do contador. A biblioteca \textit{AsconSecure} faz a mesma reserva por conta própria. Como essa gravação atrasava a transmissão dos quadros de contador múltiplo de 32, no caminho Ascon o nó cifra o quadro seguinte logo depois de cada transmissão, de modo que, no instante agendado, só o rádio trabalha; o custo é que a leitura do sensor vai ao ar com cinco segundos de idade. Os dados de aplicação (temperatura e iluminância) são, nesta etapa, sintéticos: senoides com período conhecido, o que permite verificar a decifragem no receptor por inspeção. Os metadados de recepção, que são as variáveis dos modelos, não dependem deles.

O \textit{gateway} divide o trabalho entre os dois núcleos do processador. No primeiro, uma interrupção do transceptor sinaliza cada quadro recebido; o firmware lê o pacote, registra RSSI e SNR, recoloca o rádio em escuta e só então valida o quadro. No caminho de linha de base a validação segue a ordem do LoRaWAN: verifica a janela anti-replay pelo contador em claro, autentica e decifra, e apenas depois de a autenticação conferir avança a janela, mantida em memória volátil. No caminho Ascon a biblioteca confere a tag antes do contador, de modo que um quadro forjado com contador alto não bloqueia o nó legítimo, e grava na memória não volátil o maior contador aceito de cada nó, o que faz a proteção contra repetição sobreviver a um reinício do \textit{gateway}. Nesse caminho o \textit{gateway} mede ainda, por quadro, o tempo da decifragem e o da validação completa, lidos do contador de ciclos do processador, e os envia ao \textit{backend} junto com o tempo de entrega; a linha serial de doze campos não muda. O registro JSON com os doze campos sai pela interface serial antes de qualquer envio pela rede, sob exclusão mútua para que nenhum diagnóstico o corte ao meio. O quadro é então depositado, sem espera, em uma fila para o segundo núcleo. Nesse núcleo, uma tarefa mantém a conexão Wi-Fi e a sincronização de relógio por NTP, agrupa os quadros em lotes de até vinte, envia ao \textit{backend} e, em caso de falha, guarda até 512 registros e tenta novamente com espera crescente de 1 a 30 segundos. Cada quadro entregue mais de quinze segundos após a recepção é marcado como reenviado, e o instante UTC de recepção é reconstruído a partir do relógio monotônico do dispositivo, de modo que quadros recebidos antes da sincronização de relógio saem com hora correta. A cada trinta segundos o \textit{gateway} emite um sinal de vida com os contadores de quadros recebidos, rejeitados por replay, por autenticação e por CRC, e com o estado da fila e da conexão.

O \textit{backend} é uma API em FastAPI com banco SQLite. A rota de ingestão valida cada quadro contra um esquema que rejeita campos desconhecidos, para que um nome de campo errado produza um erro imediato em vez de uma coluna perdida, e insere com um índice único sobre \textit{gateway}, sessão, dispositivo e contador, o que torna o reenvio idempotente. Rotas de consulta devolvem quadros, resumo por sessão e agregados de uma sessão, e uma página de monitoramento servida pela própria API exibe, em tempo quase real, as séries de RSSI, SNR e intervalo, com marcas nos reinícios do nó, e, para o caminho Ascon, as séries de tempo de decifragem, de validação e de entrega. Os tempos de processamento chegam em colunas próprias e, quando um quadro é inserido primeiro a partir do registro serial e depois reenviado pelo \textit{gateway}, as colunas vazias são preenchidas sem sobrescrever as já gravadas. A separação entre salto por reinício e perda de rádio, descrita na Seção~\ref{sec:resultados_infra}, é feita por um módulo único, usado tanto pela API quanto pela ferramenta de análise do registro serial, para que as duas fontes nunca discordem. O \textit{backend} tem 27 testes automatizados, executados contra um banco temporário para nunca tocar o banco das coletas.

Cinco suítes de teste acompanham o firmware e foram todas executadas na própria placa, com 65 casos no total: verificação de inicialização (17 casos: chip, memória, pinos, oscilador e parâmetros de rádio), executada em cada uma das quatro placas usadas no projeto; codificação e decodificação do quadro de linha de base com rejeição de quadros forjados e repetidos (14 casos); formatação byte a byte do registro JSON e do objeto HTTP, sessões e intervalos (16 casos); persistência do contador através de reinícios (4 casos); e, para o caminho Ascon, ida e volta do quadro em uma única placa, com adulteração de dados, tag, cabeçalho e contador, nó desconhecido, repetição e quadro forjado com contador alto (14 casos), executada em 7 de outubro de 2026, antes da primeira transmissão com Ascon, e de novo em 16 de novembro, na placa que assumiu o papel de nó. As três suítes do meio só compilavam até 17 de novembro de 2026, quando foram executadas; a do registro JSON falhou em dois casos, ambos por erro no valor esperado escrito no próprio teste (uma contagem de aspas e um instante de referência errado em 40 minutos), e nenhum no firmware. O episódio justifica a regra adotada de que uma suíte só conta como verificação depois de executada no hardware.



%%%%%%%%%%%%%%%%%%%%%%%%%%%%%%%%%%%%%%%%%%%%%%%%%%%%

